% Lettre aux enfants (deux pages). Aucune anecdote inventée.
% À PERSONNALISER PAR L'AUTEUR (to be personalised by the author) : cette lettre est un cadre ; l'auteur peut y
% ajouter ses propres souvenirs, que les agents n'ont pas le droit d'inventer.
\chapter*{Lettre aux enfants}
\addcontentsline{toc}{chapter}{Lettre aux enfants}
\markboth{Lettre aux enfants}{Lettre aux enfants}

\noindent\emph{Ilene, Alexis, Aurelia,}

\medskip
Ce livre parle de machines, de galaxies et de cristaux, mais il a été écrit en pensant à vous. Je voudrais vous dire,
en quelques lignes, les trois choses que cette année m'a apprises et que j'aimerais vous transmettre.

\paragraph{La curiosité d'abord} Tout ce qui est raconté ici a commencé par une question : « et si c'était vrai ? »,
ou « et si c'était faux ? ». Poser des questions n'est pas un défaut d'élève qui n'a pas compris ; c'est le début de
tout. Les machines de ce livre savent répondre à presque tout. Elles ne savent pas encore choisir les questions qui
méritent d'être posées. Ça, c'est votre travail, et il vous appartient. Gardez l'habitude de demander « pourquoi ? »
et « comment le sait-on ? », même aux adultes, même aux livres, même à celui-ci.

\paragraph{Vérifier son travail} Vous avez vu, dans ces pages, qu'une machine très savante peut annoncer « c'est
fini » alors que rien n'est fait. Les humains aussi, et moi le premier. La seule protection que je connaisse, c'est
de vérifier, et de faire vérifier par quelqu'un d'autre. Refaire un calcul d'une autre façon. Essayer un cas où l'on
sait que la réponse doit être fausse, pour voir si l'on s'en aperçoit. Demander à un ami de relire. Ce n'est pas
de la méfiance : c'est du respect pour ce qu'on fait, et pour ceux qui vont s'en servir.

\paragraph{Reconnaître ses erreurs} C'est la leçon la plus difficile, et la plus importante. Cette année, j'ai
publié des choses fausses. Je les ai retirées, et j'ai écrit pourquoi, au lieu de les effacer en silence. Ce n'était
pas agréable. Mais chaque erreur reconnue a rendu le travail suivant meilleur, et chaque erreur cachée l'aurait rendu
plus fragile. Une personne qui dit « je me suis trompé » n'est pas plus faible qu'une autre ; elle est plus digne de
confiance. Vous aurez le droit de vous tromper toute votre vie. Ce que je vous souhaite, c'est d'avoir aussi le
courage de le dire.

\medskip
Il y a autre chose que je voudrais vous dire, et que les machines ne savent pas faire : une découverte n'a de sens
que si elle sert à quelqu'un. J'ai essayé, dans ce livre, de choisir des questions qui comptent pour les gens, et de
dire honnêtement ce que nous savons et ce que nous ne savons pas. Si un jour vous travaillez avec des machines
encore plus fortes que celles-ci, ce qui est probable, souvenez-vous que c'est à vous de décider ce que leur travail
veut dire.

À la toute fin de ce livre, cinq images vous attendent. Je ne vous dirai pas ce qu'elles cachent ; vous le trouverez.

\bigskip
\begin{flushright}
\emph{Papa}
\end{flushright}
